\section{The Everettian diagnosis}
\label{sec:branching}

This section introduces a minimal Everettian measurement package and asks whether bare unitary branching supplies the fixation bridge isolated in Section~\ref{sec:semanticcore}. It proves that Everettian physics is non-selective in the sense fixed in Section~\ref{sec:targetclass}, separates branch-local determinacy from diachronic reference, and shows that decoherence, centered descriptions, self-location, and decision theory do not by themselves fix the referent of pre-measurement discourse.

\subsection{The Everettian base package}

The paper needs only a small amount of quantum mechanics. For the present purpose an Everettian
measurement episode can be modeled by a pre-measurement state, an entangling measurement
interaction, and a decoherence regime that renders the resulting outcome-sectors approximately
dynamically autonomous for the duration of the measurement record.

\begin{definition}[Everettian base package for a measurement episode]
\label{def:everett_base}
A measurement episode satisfies the \emph{Everettian base package} when all of the following hold:
\begin{enumerate}
    \item \textbf{Unitary evolution.} The total post-interaction state is obtained by a unitary measurement interaction $U_M$ applied to the pre-measurement state.
    \item \textbf{Outcome-sector multiplicity.} There exists a decomposition
    \[
    \Psi_M=\sum_{i\in I}\alpha_i\psi_i,
    \qquad |I|\ge 2,
    \qquad \alpha_i\neq 0 \text{ for each } i,
    \]
    where each $\psi_i$ is associated with a distinct stable record/outcome package $(\Cont_i,\Rec_i,o_i)$.
    \item \textbf{Decoherence.} The sectors $\psi_i$ are pairwise approximately non-interfering for the purposes of the measurement record, so that each supports branch-local record stability.
    \item \textbf{No selector in the bare physics.} No additional collapse postulate, hidden-variable actualization rule, or other branch-privileging physical structure is added to the bare unitary description.
\end{enumerate}
\end{definition}

The definition deliberately builds in only what an Everettian appeals to when claiming that unitary
quantum theory together with decoherence already explains measurement: global unitary evolution,
multiple nonzero outcome-sectors, and record stability without a collapse or actualization rule
\citep{everett1957,zurek2003,wallace2012}.

\begin{lemma}[Branch-local determinacy]
\label{lem:branch-local-determinacy}
Suppose a measurement episode satisfies Definition~\ref{def:everett_base}. Then each decohered
outcome-sector contains a determinate branch-local record and a determinate branch-local
observer-state.
\end{lemma}

\begin{proof}
By clause (2) of Definition~\ref{def:everett_base}, each sector $\psi_i$ is associated with a stable
record/outcome package $(\Cont_i,\Rec_i,o_i)$. By clause (3), decoherence suppresses interference
between the sectors for the purposes of the measurement record, so each such package is dynamically
stable in the branch-local sense relevant to reporting, memory, and within-branch experience. Hence
each decohered sector contains a determinate branch-local record and branch-local observer-state.
\end{proof}

\begin{lemma}[Unitary multiplicity lemma]
\label{lem:unitary_multiplicity}
Suppose a measurement episode satisfies Definition~\ref{def:everett_base}. Then the post-measurement
state contains at least two nonzero physically realized outcome-sectors.
\end{lemma}

\begin{proof}
By clause (2) of Definition~\ref{def:everett_base}, the post-measurement state has the form
\[
\Psi_M=\sum_{i\in I}\alpha_i\psi_i
\]
with $|I|\ge 2$ and $\alpha_i\neq 0$ for each $i\in I$. A nonzero coefficient means that the
corresponding sector contributes nontrivially to the physical state-description of the episode. Since
there are at least two such coefficients, the post-measurement state contains at least two nonzero
outcome-sectors. By the intended Everettian reading, those sectors are physically realized parts of the
same global state rather than merely hypothetical alternatives. Hence the episode contains at least two
nonzero physically realized outcome-sectors.
\end{proof}

\begin{lemma}[Entanglement correlates without privileging]
\label{lem:entangle_no_selector}
Let a measurement episode satisfy Definition~\ref{def:everett_base}. Then the entangling measurement
interaction correlates system, apparatus, and observer records across the distinct outcome-sectors,
but does not by itself privilege one such sector as the unique diachronic continuation relevant to
pre-measurement measurement discourse.
\end{lemma}

\begin{proof}
In the Everettian package, the measurement interaction produces a correlated state of the schematic
form
\[
\sum_{i\in I} \alpha_i\, \lvert s_i\rangle\lvert A_i\rangle\lvert O_i\rangle,
\]
where distinct system-outcome states are correlated with distinct apparatus and observer-record
states. That establishes the correlation claim: each sector carries its own determinate branch-local
triple and hence its own record/outcome package.

What must be shown is the absence of privileging. The entangling state specifies, for each $i$, what
the record is in that sector. It does not specify that one $i$ is the unique referent of pre-measurement
first-person discourse. To privilege one sector, the physical theory would need an additional rule of the
form: exactly one correlated sector is the continuation relevant to the pre-measurement subject. But
clause (4) of Definition~\ref{def:everett_base} excludes any such added selector in the bare physics.
Therefore entanglement supplies correlated multiplicity, not a unique diachronic referent.
\end{proof}

\begin{lemma}[Decoherence yields stability, not singularity]
\label{lem:decoherence_stability_only}
If a measurement episode satisfies Definition~\ref{def:everett_base}, then decoherence can explain the
stability and effective autonomy of each branch-local record without thereby selecting which branch is
the referent of pre-measurement diachronic discourse.
\end{lemma}

\begin{proof}
By clause (3) of Definition~\ref{def:everett_base}, decoherence renders the sectors approximately
non-interfering for the purposes of the measurement record. This yields the standard Everettian
achievement: each sector supports a stable branch-local record and observer-state.

However, branch-local stability is a within-sector property. It tells us that if a sector is taken as the
relevant one, its record is robust enough for local memory and reporting. It does not add any physical
fact of the form \enquote{sector $i$ rather than sector $j$ is the unique continuation relevant to
pre-measurement discourse}. Such a fact would be a selector. Since clause (4) excludes selectors in the
bare physics, decoherence yields stability but not singularity of referent.
\end{proof}

\begin{lemma}[No supervenient diachronic selector]
\label{lem:physical-parity-under-nonselection}
Let $T$ be physically non-selective with respect to a measurement process, and let $O_i$ and $O_j$
be two post-measurement observer-successors realized by $T$. Suppose that $T$ supplies no physical
law, state variable, dynamically privileged trajectory, or physically relevant asymmetry that distinguishes
$O_i$ from $O_j$ with respect to diachronic continuation. Then no unique diachronic continuation
relation privileging $O_i$ over $O_j$ supervenes on the physical facts delivered by $T$ alone.
\end{lemma}

\begin{proof}
To say that a unique diachronic continuation relation supervenes on the physical facts delivered by $T$
is to say that those facts suffice to determine, without further addition, that one successor is the
privileged continuation rather than the other. But by hypothesis, the physical description delivered by
$T$ contains no law, state variable, privileged trajectory, or other physically relevant asymmetry
distinguishing $O_i$ from $O_j$ with respect to diachronic continuation. Any relation that nevertheless
privileges $O_i$ over $O_j$ would therefore introduce an asymmetry not fixed by the physical facts
supplied by $T$. Hence no such unique privileging relation supervenes on the physical facts delivered
by $T$ alone.
\end{proof}

\begin{corollary}[Everettian non-selection]
\label{cor:everettian-nonselection}
\label{prop:no_selector_base}
In a branching Everettian measurement model satisfying physical non-selection, the physical structure
alone determines neither an explicit selector nor a supervenient unique diachronic referent for
pre-measurement first-person discourse.
\end{corollary}

\begin{proof}
By Lemma~\ref{lem:unitary_multiplicity}, the Everettian base package contains at least two nonzero
physically realized outcome-sectors, each with its own branch-local record and continuant. By
Definition~\ref{def:physical-nonselection}, the package contains no physical law, state variable, or
dynamically privileged trajectory distinguishing one such successor as physically privileged. The
absence of an explicit selector is therefore immediate. Applying Lemma~\ref{lem:physical-parity-under-nonselection}
to any two such realized successors yields the absence of a supervenient unique diachronic referent.
\end{proof}

\begin{theorem}[Everettian non-selection theorem]
\label{thm:everett_nonselection}
Every measurement episode satisfying the Everettian base package is physically non-selective in the
sense of Definition~\ref{def:physical-nonselection}.
\end{theorem}

\begin{proof}
To be physically non-selective, a measurement episode must realize multiple stable successor histories
while supplying no physical law, state variable, or dynamically privileged trajectory that singles out
one such history as physically distinguished. The first condition follows from Lemma~\ref{lem:unitary_multiplicity};
branch-local determinacy of the realized histories follows from Lemma~\ref{lem:branch-local-determinacy};
and the absence of a physically distinguishing selector is built into clause (4) of
Definition~\ref{def:everett_base} and then made explicit by Corollary~\ref{cor:everettian-nonselection}.
Hence every measurement episode satisfying the Everettian base package is physically non-selective.
\end{proof}
\begin{figure}[htbp]
\centering
\begin{tikzpicture}[
    node distance=10mm and 10mm,
    every node/.style={align=center},
    box/.style={draw=accent, rounded corners=2pt, fill=softaccent, inner sep=5pt, text width=0.22\textwidth},
    arrow/.style={-Latex, thick, draw=accent}
]
\node[box] (pre) {Pre-measurement subject $O_0$\\single anticipatory discourse\\single evidential target};
\node[box, below left=of pre] (up) {Branch $b_{\uparrow}$\\stable record $\Rec_{\uparrow}$\\continuant $\Cont_{\uparrow}$};
\node[box, below right=of pre] (down) {Branch $b_{\downarrow}$\\stable record $\Rec_{\downarrow}$\\continuant $\Cont_{\downarrow}$};
\draw[arrow] (pre) -- node[left,midway,font=\small]{unitary evolution + decoherence} (up);
\draw[arrow] (pre) -- node[right,midway,font=\small]{unitary evolution + decoherence} (down);
\node[below=8mm of $(up)!0.5!(down)$, text width=0.68\textwidth] {The physical package yields multiplicity, correlation, and branch-local stability. What it still does \emph{not} yield is a selector fixing which post-measurement history is the referent of the pre-measurement subject's diachronic discourse.};
\end{tikzpicture}
\caption{Minimal branching diagram for the theorem ladder. The Everettian base package yields multiple stable post-measurement records, but it does not physically select which history is the diachronic referent of the pre-measurement discourse.}
\label{fig:minimal-branching}
\end{figure}

\subsection{Branch-relative determinacy versus diachronic reference}

Once the Everettian base package is in place, one may grant a great deal. Each realized branch can
contain a determinate local observer, a determinate local record, and a determinate local outcome.
What still does not follow is that the pre-measurement subject's diachronic discourse has a uniquely
fixed referent.

\begin{definition}[Successor-qualified history]
A physically realized post-measurement history $b\in\B$ is \emph{successor-qualified} for the
purposes of a given piece of diachronic discourse if it contains both
\begin{enumerate}
    \item a determinate measurement record $\Rec(b)$ appropriate to that discourse, and
    \item a post-measurement continuant $\Cont(b)$ satisfying the continuity conditions presupposed by
    the discourse between the pre-measurement subject and the future subject.
\end{enumerate}
\end{definition}

\begin{lemma}[Multiplicity of successor-qualified histories yields multiplicity of admissible assignments]
\label{lem:successorqualified}
Let $\mathcal{M}$ be a physically non-selective measurement episode and let $d$ be a record-sensitive
claim in its target discourse. If there exist distinct successor-qualified histories $b_i,b_j\in\B$ for $d$,
then $\ARef(\mathcal{M})$ contains at least two distinct admissible reference assignments relevant to
$d$.
\end{lemma}

\begin{proof}
Because $b_i$ is successor-qualified for the discourse involving $d$, it satisfies the conditions that
make the pair $(b_i,\Rec(b_i))$ eligible to serve as a referent for $d$. So there is an admissible
assignment $\rho_i$ sending $d$ to $(b_i,\Rec(b_i))$. The same reasoning applies to $b_j$, yielding an
admissible assignment $\rho_j$ sending $d$ to $(b_j,\Rec(b_j))$. Because $b_i\neq b_j$, the
assignments are distinct. Therefore there are at least two admissible assignments relevant to $d$.
\end{proof}

\begin{lemma}[Branch-relative determinacy does not entail diachronic determinacy]
\label{thm:branchlocal-not-diachronic}
\label{lem:branchrelative-not-diachronic}
The existence of determinate branch-relative post-measurement experiences does not determine which such branch-relative experience is the referent of pre-measurement diachronic first-person discourse.
\end{lemma}

\begin{proof}
Consider a branching model in which two post-measurement successors $O_1$ and $O_2$ each possess determinate branch-local experiences and stable branch-relative records. Now consider two admissible assignments: one taking $O_1$ to be the referent of the pre-measurement subject's diachronic discourse, the other taking $O_2$ to be that referent. The branch-local determinacy of each successor remains unchanged across the two assignments. What changes is only the diachronic successor-reference relation. Hence branch-local determinacy does not entail diachronic referential determinacy.
\end{proof}

\begin{corollary}[Determinacy within branches is compatible with global referential underdetermination]
A theory may deliver determinate experiences and records inside every realized branch while still
failing to determine which branch-history is the referent of pre-measurement memory, anticipation,
and evidential discourse.
\end{corollary}

\begin{proof}
This is an immediate restatement of
Theorem~\ref{thm:branchlocal-not-diachronic}. The theorem establishes that local determinacy in
 each branch and diachronic determinacy across branches are distinct conditions. The first may hold
while the second fails.
\end{proof}

\begin{proposition}[Successor multiplicity marks a structural distinction]
The difference between branch-local determinacy and diachronic referential determinacy marks a structural difference between two explanatory questions:
\begin{enumerate}
    \item what records and observer-stages exist in each realized history; and
    \item which such history is the referent of pre-measurement first-person and evidential discourse.
\end{enumerate}
\end{proposition}

\begin{proof}
The first question is answered by the branch-local facts contained in each history: what outcome is
recorded there, what observer-stage exists there, and what local memory state exists there. The second
question is answered only if the theory determines which realized history is relevant to the
pre-measurement subject's diachronic claims. A theory may answer the first question for every branch
and still leave the second unanswered. That possibility is exactly what
Theorem~\ref{thm:branchlocal-not-diachronic} shows. Therefore the distinction is structural.
\end{proof}

\subsection{Decoherence and remaining indeterminacy}

All discussion of evidential and experiential reference in this section is conditional on the target
explanatory burden fixed in Section~\ref{sec:targetclass}. The issue is not whether decoherence is
useful or real. The issue is whether decoherence, by itself, fixes diachronic reference.

\begin{definition}[Record stability under decoherence]
A measurement episode $\mathcal{M}$ satisfies \emph{decoherence-based record stability} if for every
physically realized history $b\in\B$, the record $\Rec(b)$ is dynamically stable enough to support
branch-local memory and observational reliability, and the corresponding branch remains
approximately non-interfering with the others for the purposes of the episode under discussion.
\end{definition}

\begin{proposition}[Decoherence stabilizes records without selecting referents]
\label{prop:decoherence}
Let $\mathcal{M}$ be a physically non-selective measurement episode satisfying decoherence-based
record stability. Then decoherence stabilizes the records contained in each $b\in\B$ but does not by
itself privilege one $b$ as the unique referent of pre-measurement diachronic discourse.
\end{proposition}

\begin{proof}
By decoherence-based record stability, each history $b\in\B$ contains a dynamically stable record
$\Rec(b)$ that may support branch-local memory and branch-local observational discourse. This is the
positive achievement of decoherence.

To show that decoherence does not by itself privilege a unique diachronic referent, suppose for
contradiction that it does. Then decoherence would provide a selector: it would reduce the admissible
reference possibilities relevant to the discourse to exactly one privileged assignment. But
Decoherence-based record stability, as defined here, supplies only local stability and effective
non-interference of the realized histories. It says that each realized branch carries a stable record; it
does not say that one stable branch is physically privileged over the others as \emph{the}
continuation relevant to pre-measurement discourse. Adding such a privilege would amount to adding
a selector not already contained in the bare statement of decoherence-based stability. Therefore
decoherence stabilizes records without by itself selecting referents.
\end{proof}

\begin{theorem}[Decoherence non-selection theorem]
\label{thm:decoherence_nonselection}
If a theory's post-measurement contribution to a measurement episode consists only of the physical
realization of multiple branch-histories together with decoherence-based record stability, then the
theory does not by that contribution alone fix diachronic reference for measurement records.
\end{theorem}

\begin{proof}
Assume the theory's post-measurement contribution consists only of the physical realization of
multiple histories together with decoherence-based record stability. By
Proposition~\ref{prop:decoherence}, that contribution stabilizes the records in each realized history but
does not itself provide a selector. If the episode is physically non-selective, then at least two
admissible reference assignments remain available. Therefore the theory's post-measurement
contribution, as specified, does not fix diachronic reference for the measurement records.
\end{proof}

\subsection{Centered descriptions, self-location, and decision theory}

The most common way of resisting the previous result is to say that Everettian theory has resources for
reasoning under multiplicity: centered-world semantics, self-locating uncertainty, and
decision-theoretic branch weighting \citep{lewis1979,greaves2004,sebenscarroll2018}. The present claim is not that such resources are incoherent. The claim is that they presuppose a prior typing of successors.

\begin{definition}[Centered continuation description]
A \emph{centered continuation description} for a measurement episode $\mathcal{M}$ is a pair
$(b,c_b)$ where $b\in\B$ is a physically realized post-measurement history and $c_b$ is a center-marker
intended to treat the post-measurement continuant in $b$ as a candidate continuation of the
pre-measurement subject for first-person or evidential discourse.
\end{definition}

\begin{definition}[Centered enrichment]
A \emph{centered enrichment} of a physically non-selective measurement episode $\mathcal{M}$ is the
passage from the uncentered history family $\B$ to a family of centered continuation descriptions
\[
\mathcal{C}_M=\{(b,c_b): b\in\B\}.
\]
The enrichment is \emph{selector-free} if it adds no rule that privileges exactly one center $(b,c_b)$ as
the referent of pre-measurement diachronic discourse.
\end{definition}

\begin{lemma}[Centered enrichment preserves multiplicity]
\label{lem:centered_preserves}
Let $\mathcal{M}$ be a physically non-selective measurement episode and let $\mathcal{C}_M$ be a
selector-free centered enrichment of $\mathcal{M}$. Then if $|\B|\ge 2$, the centered enrichment
contains at least two distinct candidate centers.
\end{lemma}

\begin{proof}
Because $\mathcal{M}$ is physically non-selective, $|\B|\ge 2$. So there exist distinct histories
$b_i\neq b_j$ in $\B$. By definition of centered enrichment, both give rise to centered continuation
descriptions $(b_i,c_{b_i})$ and $(b_j,c_{b_j})$. Since the underlying histories are distinct, these
centered descriptions are distinct as well. Therefore the centered enrichment contains at least two
distinct candidate centers.
\end{proof}

\begin{proposition}[Centered enrichment does not by itself select a referent]
\label{prop:centered_not_enough}
A selector-free centered enrichment of a physically non-selective measurement episode does not by
itself determine which centered continuation description is the referent of pre-measurement diachronic
measurement discourse.
\end{proposition}

\begin{proof}
By Lemma~\ref{lem:centered_preserves}, selector-free centered enrichment preserves multiplicity of
candidate centers. To determine which centered continuation description is the referent of
pre-measurement discourse, one would need a rule that singles out one center from the enriched set
$\mathcal{C}_M$. But selector-free centered enrichment, by definition, adds no such rule. So the
centered enrichment redescribes the multiplicity of candidate continuants; it does not eliminate that
multiplicity or privilege one member of it as the unique referent.
\end{proof}

\begin{definition}[Typed successor set]
A set $\Sigma_M$ of post-measurement continuants is a \emph{typed successor set} for a piece of
diachronic measurement discourse iff each member of $\Sigma_M$ is already fixed as an admissible
candidate continuation of the pre-measurement subject for the truth conditions of that discourse.
\end{definition}

\begin{lemma}[Self-locating uncertainty requires a typed successor set]
\label{lem:selfloc_typed}
A self-locating uncertainty rule can guide first-person credence only if it ranges over a typed
successor set.
\end{lemma}

\begin{proof}
A self-locating uncertainty rule distributes credence across candidate future centers or successors. For
that distribution to be first-personally meaningful, the rule's domain must already consist of the
candidates relevant to the subject whose uncertainty is in question. If the domain were not already
typed in that way, the rule would distribute credence across a merely physical collection of
post-measurement continuants without having fixed that these are the continuants relevant to the
present subject's anticipatory discourse. In that case the rule would fail to determine the truth
conditions of first-person claims such as \enquote{I will observe $o$}. So self-locating uncertainty
requires a typed successor set.
\end{proof}

\begin{definition}[Decision-theoretic branching rule]
A \emph{decision-theoretic branching rule} for a measurement episode $\mathcal{M}$ is any rule that
assigns utilities, preferences, or credences across a set of candidate post-measurement continuants or
records and is then used to guide pre-measurement deliberation.
\end{definition}

\begin{lemma}[Decision rules require typed successor events]
\label{lem:typed_successors}
A decision-theoretic branching rule can guide action only if the successor events or records to which it
assigns utilities or credences are already typed as the relevant successors for the discourse being used by
the agent.
\end{lemma}

\begin{proof}
A decision-theoretic rule guides action only if its outputs bear on what the present agent should expect,
prefer, or do. To have that bearing, the successor events or records over which the rule ranges must
already be the successor events or records relevant to the present agent. If those items are not already
typed as the relevant successors for the discourse in which the decision is being made, then the rule's
outputs are not yet action-guiding in the intended first-person sense. The rule would at best distribute
weights over an untyped collection of physical continuants without thereby determining which of them
are the future experiences or records to which the present agent's claims refer. Therefore the decision
rule requires successor events or records that are already referentially typed for the discourse in
question.
\end{proof}

\begin{proposition}[Weights over branches are not selectors]
\label{prop:weightsnotselectors}
Assigning amplitudes, credences, utilities, or caring measures to realized branches does not by itself
fix which history is the referent of pre-measurement diachronic discourse.
\end{proposition}

\begin{proof}
A selector determines which realized history is the relevant referent. A weight assignment instead
orders or measures the candidate histories while leaving the domain of candidates in place. Unless one
already knows what the weighted items are the candidates \emph{for}, the weights do not fix which
history the pre-measurement discourse is about. They therefore presuppose rather than provide the
needed referential typing.
\end{proof}

\begin{theorem}[Everettian downstream theorem]
\label{thm:downstream}
Centered descriptions, self-locating uncertainty, and decision-theoretic branch-weighting are
downstream of the Everettian non-selection problem. They presuppose successor typing and therefore
cannot by themselves supply the diachronic reference fixation whose absence constitutes the problem
identified in this paper.
\end{theorem}

\begin{proof}
By Proposition~\ref{prop:centered_not_enough}, centered enrichment without a selector preserves
rather than eliminates multiplicity of candidate continuants. By Lemma~\ref{lem:selfloc_typed}, any
self-locating treatment can guide first-person credence only if its domain is already a typed successor
set. By Lemma~\ref{lem:typed_successors}, a decision-theoretic branching rule can guide action only if
its successor events are already typed as the relevant successors for the present agent's discourse. But
Theorem~\ref{thm:everett_nonselection} and
Theorem~\ref{thm:branchlocal-not-diachronic} show that the bare Everettian physical package leaves
multiple realized histories in play without fixing which one is the relevant diachronic referent.
Therefore centered descriptions, self-location, and decision theory begin only after the missing
referential work has been done. They are downstream of the non-selection problem rather than solutions to it.
\end{proof}

Section~\ref{sec:branching} has established that the standard Everettian physical package yields multiple stable branch-local records and no unique diachronic referent, and that self-location and decision-theoretic replies presuppose rather than remove that problem. The next section turns to how measurement records function as evidential anchors and what happens to theory-level confirmation when admissible reference assignments diverge.
